\section{Methodology}
\label{sec:methodology}

\subsection{Overview of the Cognitive Research Memory Engine (CRME)}

The Cognitive Research Memory Engine (CRME) is a structured research memory framework designed to support continuity, organization, traceability, and knowledge evolution throughout the research lifecycle.

Unlike conventional documentation-based approaches that primarily store research outputs as isolated files, CRME models research knowledge as a collection of interconnected and evolving research objects. Each research object maintains semantic information, contextual state, relationships, temporal information, and provenance records.

The primary objective of CRME is to provide a unified infrastructure for managing research context, historical interactions, project evolution, semantic relationships, and generated artifacts. The framework integrates temporal memory management, project-level organization, knowledge representation, and provenance tracking within a modular architecture.

As illustrated in Fig.~\ref{fig:crme_architecture}, CRME follows a layered architecture consisting of the Research Interaction Layer, Cognitive Coordinator, Specialized Memory Engines, Research Memory Repository, and Reasoning and Intelligence Layer.

This architecture enables continuous interaction between researchers and accumulated research knowledge, allowing previous research states to be recovered and reused during future activities.


\subsection{CRME Architecture}

The overall architecture of CRME is presented in Fig.~\ref{fig:crme_architecture}. The framework is organized around a Cognitive Coordinator that manages communication between research activities and memory components.

The Cognitive Coordinator acts as the central orchestration mechanism responsible for context management, memory coordination, and knowledge routing.

Unlike passive storage systems, the coordinator enables dynamic creation, update, and retrieval of research knowledge.


\subsubsection{Session Engine}

The Session Engine provides temporal continuity across research interactions by capturing session states and maintaining historical snapshots.

Each research session is represented as:

\[
S_t =
(G_t,D_t,M_t,A_t)
\]

where $G_t$ represents research goals, $D_t$ represents decisions, $M_t$ represents interaction messages and contextual information, and $A_t$ represents generated artifacts.

Through session-based memory representation, CRME preserves the evolution of research activities across multiple interactions.


\subsubsection{Project Engine}

The Project Engine provides research lifecycle management capabilities by representing project-level structures, including objectives, milestones, dependencies, and progress information.

By connecting research activities with project-level states, CRME provides a higher-level representation of research development.


\subsubsection{Knowledge Graph Engine}

The Knowledge Graph Engine represents semantic relationships among research objects.

The research memory graph is modeled as:

\[
G=(V,E)
\]

where $V$ represents research objects and $E$ represents semantic relationships between objects.

This graph-based representation enables structured knowledge organization and supports retrieval and future reasoning operations.


\subsubsection{Research Memory Repository}

The Research Memory Repository acts as the central storage layer for structured research information.

It maintains research objects, metadata, session histories, decision records, and generated artifacts.

Unlike traditional repositories that mainly focus on document storage, CRME preserves contextual relationships associated with each research entity.


\subsubsection{Reasoning and Intelligence Layer}

The Reasoning and Intelligence Layer provides mechanisms for knowledge retrieval, contextual analysis, and future reasoning extensions over accumulated research memory.

The feedback loop illustrated in Fig.~\ref{fig:crme_architecture} enables extracted knowledge to influence future research activities through the Cognitive Coordinator.


\subsection{Research Memory Lifecycle}

Research activities are iterative processes in which ideas, hypotheses, experiments, and conclusions continuously evolve.

As illustrated in Fig.~\ref{fig:research_lifecycle}, CRME models this evolution through a research memory lifecycle consisting of idea formation, hypothesis development, planning, experimentation, artifact generation, and knowledge feedback.

This lifecycle-oriented design distinguishes CRME from static documentation systems by representing research as an evolving knowledge process rather than a collection of independent artifacts.


\subsection{Research Object Representation}

A fundamental design principle of CRME is representing research knowledge through structured Research Objects.

A Research Object is formally defined as:

\[
RO_i =
(ID_i,Type_i,Meta_i,Rel_i,Temp_i,Prov_i,Art_i)
\]

where:

\begin{itemize}
\item $ID_i$ represents the unique identifier.
\item $Type_i$ defines the research object category.
\item $Meta_i$ represents descriptive attributes.
\item $Rel_i$ captures semantic relationships.
\item $Temp_i$ maintains temporal information.
\item $Prov_i$ records provenance information.
\item $Art_i$ represents associated artifacts.
\end{itemize}

The object type space includes:

\[
Type =
\{Idea,Hypothesis,Decision,Experiment,Result,Artifact,Publication\}
\]

This representation enables CRME to preserve both research outcomes and the historical context that produced them.


\subsection{Provenance-Aware Research Memory}

CRME incorporates provenance-aware memory management to improve traceability between decisions, processes, and research outcomes.

For each research entity, provenance information is represented as:

\[
Prov_i =
(Source_i,Activity_i,Agent_i,Time_i)
\]

This mechanism enables researchers to investigate how artifacts were generated, which decisions influenced their creation, and how previous research states contributed to current outcomes.


\subsection{Implementation Realization}

The proposed methodology has been implemented using a modular Python-based architecture.

The mapping between conceptual components and implementation modules is summarized in Table~\ref{tab:implementation_mapping}.

\begin{table}[h]
\centering
\caption{Mapping Between CRME Components and Implementation Modules}
\label{tab:implementation_mapping}

\begin{tabular}{ll}
\hline
Conceptual Component & Implementation Module\\
\hline
Memory Management & memory\_engine.py\\
Session Continuity & session\_engine.py\\
Project Management & project\_engine.py\\
Knowledge Relationships & graph\_engine.py\\
Semantic Retrieval & semantic\_engine.py\\
Evaluation Framework & evaluation/ modules\\
\hline
\end{tabular}

\end{table}


\subsection{Evaluation Integration}

CRME evaluation is integrated into the proposed methodology through benchmark comparison, ablation analysis, and scalability evaluation.

The evaluation framework measures:

\begin{itemize}
\item Knowledge Organization Score (KOS),
\item Decision Traceability Score (DTS),
\item Research State Retention (RSR),
\item Cognitive Retrieval Time (CRT).
\end{itemize}


\subsection{Summary}

CRME introduces a structured methodology for managing evolving research knowledge through coordinated memory components, semantic representation, temporal continuity, and provenance-aware tracking.

By integrating research objects, session memory, project evolution, semantic relationships, and artifact lineage, CRME transforms research memory from passive storage into an evolving knowledge infrastructure.

